% !TEX root = ../main.tex

\chapter{系统总体设计}


\section{系统角色模型}


本文系统采用 Alice - ManagerA - Notary - ManagerB - Bob 的跨链协作模型。该模型将跨链交易中的资金流、消息流和控制流拆分为不同角色，从而提高各环节的可分析性。


Alice 是跨链交易发起方。她负责生成 secret，计算 hashLock，并在源链发起跨链请求。Alice 在源链上锁定资产或中间代币后，等待目标链完成对应资产释放。Bob 是目标链上的接收方，最终应在目标链收到对应资产。在最简单的跨链场景中，Alice 和 Bob 可以是两个不同用户，也可以是同一用户在不同链上的地址。


ManagerA 和 ManagerB 是两侧资产池或跨链 TokenPool 的管理者。ManagerA 位于源链一侧，负责接收、兑换或锁定 Alice 的资产；ManagerB 位于目标链一侧，负责为 Bob 创建对应锁定并在满足条件时释放资产。Manager 并不是任意转移用户资金的中心化操作者，其行为受到合约状态、hashLock、timeLock 和权限控制约束。


Notary 是公证与中继节点，负责监听源链事件并协调目标链操作。当 Alice 在源链发起跨链请求后，源链合约会产生事件，Notary 从事件中读取 swapId、hashLock、amount、recipient 等参数，并将其转发给目标链 ManagerB。Notary 的角色重点在于跨链消息传递和流程协调，而不是最终托管资产。


这种角色模型的优点在于结构清晰。Alice 和 Bob 是用户端，Manager 处理资金池，Notary 处理消息协调，HTLC 或 adaptor signature 处理条件锁定。论文后续章节中的安全性和兼容性分析也都围绕这一角色模型展开。


\section{系统模块划分}


从工程实现角度看，系统可以划分为智能合约模块、后端协调模块、scriptless 适配模块和前端可视化模块。系统的整体分层架构如图~\ref{fig:crosschain_arch} 所示，自上而下分为前端可视化层、后端协调层和链上机制层，五类角色分别落在不同层次：用户端（Alice/Bob）通过前端层发起和观察交易，Notary 位于后端协调层负责跨链消息中继，ManagerA/ManagerB 与具体的链上机制（HTLC 或 adaptor signature）位于链上机制层。

\begin{figure}[htbp]
  \centering
  % 跨链系统分层架构图
\begin{tikzpicture}[
    font=\footnotesize,
    layer/.style={rectangle, draw, rounded corners, minimum height=3.2em, align=center, fill=black!3},
    mod/.style={rectangle, draw, minimum height=2.2em, minimum width=7.5em, align=center, fill=white},
    arrow/.style={-{LaTeX}, thick},
    biarrow/.style={{LaTeX}-{LaTeX}, thick},
  ]

  % ===== 前端可视化层 =====
  \node[mod] (fe1) {跨链交换页面\\（产品原型）};
  \node[mod, right=1.2cm of fe1] (fe2) {HTLC 演示页面\\（流程可视化）};
  \begin{scope}[on background layer]
    \node[layer, fit=(fe1)(fe2), inner sep=1.4ex, label={[anchor=west,xshift=-0.2cm]west:\textbf{前端可视化层}}] (felayer) {};
  \end{scope}

  % ===== 后端协调层 =====
  \node[mod, below=2.4cm of fe1, xshift=-0.6cm] (be1) {API 路由层};
  \node[mod, right=0.6cm of be1] (be2) {Coordinator\\分发层};
  \node[mod, right=0.6cm of be2] (be3) {各链适配模块\\(ETH/BSC/Conflux/\\Fabric/BTC)};
  \begin{scope}[on background layer]
    \node[layer, fit=(be1)(be2)(be3), inner sep=1.4ex, label={[anchor=west,xshift=-0.2cm]west:\textbf{后端协调层}}] (belayer) {};
  \end{scope}

  % ===== 链上机制层 =====
  \node[mod, below=2.4cm of be1, xshift=0.2cm] (cl1) {智能合约链\\HTLC 状态机\\(CrossChainTokenPool)};
  \node[mod, right=1.6cm of cl1] (cl2) {无脚本链\\adaptor signature\\(scriptless 适配)};
  \begin{scope}[on background layer]
    \node[layer, fit=(cl1)(cl2), inner sep=1.4ex, label={[anchor=west,xshift=-0.2cm]west:\textbf{链上机制层}}] (cllayer) {};
  \end{scope}

  % ===== 层间连接 =====
  \draw[biarrow] (felayer) -- node[right]{HTTP API} (belayer);
  \draw[biarrow] (belayer) -- node[right]{事件监听 / 交易构造} (cllayer);

  % ===== 角色标注（右侧） =====
  \node[align=left, right=0.8cm of felayer.east, anchor=west] (rA) {Alice / Bob\\（用户端）};
  \node[align=left, right=0.8cm of belayer.east, anchor=west] (rN) {Notary\\（公证中继）};
  \node[align=left, right=0.8cm of cllayer.east, anchor=west] (rM) {ManagerA / ManagerB\\（资金池）};

  \draw[dashed, gray] (felayer.east) -- (rA.west);
  \draw[dashed, gray] (belayer.east) -- (rN.west);
  \draw[dashed, gray] (cllayer.east) -- (rM.west);

\end{tikzpicture}

  \caption{跨链系统分层架构}
  \label{fig:crosschain_arch}
\end{figure}


智能合约模块主要对应 CrossChainTokenPool 合约。该模块实现跨链资产池、原生币与中间代币兑换、跨链请求发起、Manager 响应、资产领取和超时退款等功能。合约中保存 CrossChainSwap 和 ManagerLock 两类状态，分别表示用户侧锁定和 Manager 侧锁定。二者通过 hashLock 和 swapId 建立关联。


后端协调模块主要由 Python 服务和跨链 coordinator 组成。该模块负责提供 API、连接不同链节点、监听合约事件、处理跨链状态、构造目标链交易并协调 Manager 响应。后端并不替代链上合约的安全判断，而是将不同链之间的消息传递和交易调用组织起来。


scriptless 适配模块主要用于展示 Schnorr adaptor signature 机制，面向 Monero、ZCash 屏蔽交易等无脚本链的跨链场景。该模块通过模拟钱包、区块链和签名流程，说明如何创建 adaptor signature、如何补全签名、如何从完整签名中提取 secret，以及如何将该机制用于跨链条件交换。由于无脚本链无法在链上表达 HTLC 的条件分支，adaptor signature 将条件锁定放入签名补全过程，使这类链仅凭签名能力即可参与原子交换。原型实现以 BTC 作为该场景的替代实现，仅用于验证机制正确性。


前端可视化模块用于展示系统流程。它包括跨链交换入口原型和 HTLC 流程演示页面。演示页面能够展示 Alice、ManagerA、Notary、ManagerB、Bob 等角色，展示 secret、hashLock、timeLock、swapId、余额变化和交易阶段，有助于理解跨链协议的执行过程。


\section{总体跨链流程}


系统总体流程可以分为正常完成路径和超时退款路径。


在正常完成路径中，Alice 首先生成 secret，并计算 hashLock = H(secret)。随后 Alice 在源链发起跨链请求，提交 hashLock、目标接收地址、金额和 timeLock 等参数。源链 ManagerA 或 TokenPool 合约将 Alice 的资产锁定，并产生 CrossChainInitiated 事件。Notary 监听到该事件后，读取跨链参数，并通知目标链 ManagerB 创建对应的 ManagerLock。ManagerB 在目标链锁定等值资产，等待 Bob 或相关执行方提交 secret。当 secret 被公开后，目标链合约验证 H(secret) 是否等于 hashLock。验证通过后，目标链释放资产给 Bob。由于 secret 已公开，源链也可以用同一 secret 完成结算。至此，一笔跨链交换完成。该路径中各角色之间的消息与资产流转顺序如图~\ref{fig:crosschain_seq_complete} 所示。

\begin{figure}[htbp]
  \centering
  % 跨链正常完成路径时序图
\begin{tikzpicture}[
    font=\footnotesize,
    actor/.style={rectangle, draw, rounded corners, minimum height=2em, minimum width=4.6em, align=center, fill=black!3},
    msg/.style={-{LaTeX}, thick},
    ret/.style={-{LaTeX}, thick, dashed},
    note/.style={font=\scriptsize\itshape, text=black!60},
  ]

  % ===== 角色生命线 X 坐标 =====
  \def\xA{0}    % Alice
  \def\xMA{2.6} % ManagerA
  \def\xN{5.2}  % Notary
  \def\xMB{7.8} % ManagerB
  \def\xB{10.4} % Bob

  % ===== 顶部角色框 =====
  \node[actor] (A) at (\xA,0) {Alice};
  \node[actor] (MA) at (\xMA,0) {ManagerA};
  \node[actor] (N) at (\xN,0) {Notary};
  \node[actor] (MB) at (\xMB,0) {ManagerB};
  \node[actor] (B) at (\xB,0) {Bob};

  % ===== 生命线（竖虚线） =====
  \def\ybot{-9.2}
  \foreach \x in {\xA,\xMA,\xN,\xMB,\xB}{
    \draw[gray, dashed] (\x,-0.45) -- (\x,\ybot);
  }

  % ===== 消息（自上而下） =====
  % 1. Alice 生成 secret / hashLock
  \node[note, right] at (\xA+0.15,-1.0) {生成 secret，计算 hashLock = H(secret)};

  % 2. Alice -> ManagerA: initiateCrossChain
  \draw[msg] (\xA,-1.6) -- node[above, font=\scriptsize]{initiate(hashLock, recipient, amount, timeLock)} (\xMA,-1.6);
  \node[note, right] at (\xMA-2.3,-2.05) {源链锁定 Alice 资产，记录 CrossChainSwap};

  % 3. ManagerA --> Notary: 触发事件
  \draw[msg] (\xMA,-2.6) -- node[above, font=\scriptsize]{CrossChainInitiated 事件} (\xN,-2.6);
  \node[note, right] at (\xN-2.0,-3.05) {解析 swapId/hashLock/amount/recipient};

  % 4. Notary -> ManagerB: 通知创建锁定
  \draw[msg] (\xN,-3.6) -- node[above, font=\scriptsize]{managerRespond(hashLock, ...)} (\xMB,-3.6);
  \node[note, right] at (\xMB-2.1,-4.05) {目标链创建 ManagerLock（较短 timeLock）};

  % 5. ManagerB -> Bob: 锁定就绪
  \draw[msg] (\xMB,-4.6) -- node[above, font=\scriptsize]{ManagerLockCreated} (\xB,-4.6);

  % 6. Bob -> ManagerB: 揭示 secret 领取
  \draw[msg] (\xB,-5.6) -- node[above, font=\scriptsize]{completeManagerLock(secret)} (\xMB,-5.6);
  \node[note, right] at (\xMB-2.4,-6.05) {验证 H(secret)=hashLock，释放资产给 Bob};

  % 7. secret 公开上链，回传至源链侧（Bob/ManagerB -> Alice）
  \draw[ret] (\xMB,-6.6) -- node[above, font=\scriptsize]{secret 已在链上公开} (\xA,-6.6);

  % 8. Alice -> ManagerA: 用 secret 完成源链结算
  \draw[msg] (\xA,-7.6) -- node[above, font=\scriptsize]{completeCrossChain(swapId, secret)} (\xMA,-7.6);
  \node[note, right] at (\xMA-2.0,-8.05) {源链结算完成，两侧状态均 completed};

  % ===== 底部箭头收尾 =====
  \foreach \x in {\xA,\xMA,\xN,\xMB,\xB}{
    \draw[gray, dashed, -{LaTeX}] (\x,\ybot+0.2) -- (\x,\ybot);
  }

\end{tikzpicture}

  \caption{正常完成路径时序图}
  \label{fig:crosschain_seq_complete}
\end{figure}


在超时退款路径中，如果 Alice 没有公开 secret，或者 Notary 未能及时协调目标链响应，或者 Bob 没有在规定时间内领取资产，则系统进入失败处理。timeLock 到期后，源链上的 Alice 可以调用退款逻辑取回被锁定资产；目标链 ManagerB 也可以对未完成的 ManagerLock 进行退款或回收。这样即使跨链流程中断，资产也不会永久停留在锁定状态。该路径如图~\ref{fig:crosschain_seq_refund} 所示，需要特别注意退款的时间顺序：目标链 ManagerLock 的 timeLock 设置得比源链更短，因此目标链先到期退款、源链后到期退款，这一不对称约束在第五章的原子性分析中有进一步论证。

\begin{figure}[htbp]
  \centering
  % 跨链超时退款路径时序图
\begin{tikzpicture}[
    font=\footnotesize,
    actor/.style={rectangle, draw, rounded corners, minimum height=2em, minimum width=4.6em, align=center, fill=black!3},
    msg/.style={-{LaTeX}, thick},
    xmsg/.style={-{LaTeX}, thick, red!60!black},
    note/.style={font=\scriptsize\itshape, text=black!60},
  ]

  % ===== 角色生命线 X 坐标 =====
  \def\xA{0}
  \def\xMA{2.6}
  \def\xN{5.2}
  \def\xMB{7.8}
  \def\xB{10.4}

  % ===== 顶部角色框 =====
  \node[actor] (A) at (\xA,0) {Alice};
  \node[actor] (MA) at (\xMA,0) {ManagerA};
  \node[actor] (N) at (\xN,0) {Notary};
  \node[actor] (MB) at (\xMB,0) {ManagerB};
  \node[actor] (B) at (\xB,0) {Bob};

  % ===== 生命线 =====
  \def\ybot{-8.6}
  \foreach \x in {\xA,\xMA,\xN,\xMB,\xB}{
    \draw[gray, dashed] (\x,-0.45) -- (\x,\ybot);
  }

  % 1. Alice 发起锁定（与正常路径相同的起点）
  \draw[msg] (\xA,-1.4) -- node[above, font=\scriptsize]{initiate(...)} (\xMA,-1.4);
  \node[note, right] at (\xMA-1.6,-1.85) {源链锁定，timeLock\_source（较长）};

  % 2. Notary 转发，ManagerB 创建锁定
  \draw[msg] (\xMA,-2.4) -- node[above, font=\scriptsize]{event} (\xN,-2.4);
  \draw[msg] (\xN,-2.9) -- node[above, font=\scriptsize]{managerRespond(...)} (\xMB,-2.9);
  \node[note, right] at (\xMB-2.2,-3.35) {目标链锁定，timeLock\_target（较短）};

  % 3. secret 始终未公开
  \node[note, red!60!black] at (\xN,-4.2) {secret 始终未公开（参与方不配合 / Notary 失效 / Bob 未领取）};
  \draw[red!50, dashed] (\xA-0.6,-4.55) -- (\xB+0.6,-4.55);

  % 4. 目标链 timeLock 先到期 -> ManagerB 退款
  \node[note, right] at (\xMB-2.6,-5.3) {timeLock\_target 先到期};
  \draw[xmsg] (\xMB,-5.8) .. controls (\xMB+0.9,-5.4) and (\xMB+0.9,-6.2) .. node[right, font=\scriptsize]{refundManagerLock} (\xMB,-6.2);

  % 5. 源链 timeLock 后到期 -> Alice 退款
  \node[note, right] at (\xA+0.15,-7.0) {timeLock\_source 后到期};
  \draw[xmsg] (\xA,-7.5) .. controls (\xA+0.9,-7.1) and (\xA+0.9,-7.9) .. node[right, font=\scriptsize]{refundCrossChain} (\xA,-7.9);

  % 6. 不对称约束说明
  \node[note, align=center, text=red!60!black] at (\xN,-8.2) {退款顺序：目标链先、源链后（timeLock\_source \textgreater\ timeLock\_target），保证安全退出};

\end{tikzpicture}

  \caption{超时退款路径时序图}
  \label{fig:crosschain_seq_refund}
\end{figure}


系统中 swapId、hashLock 和 timeLock 分别承担不同作用。swapId 是跨链请求编号，用于区分不同交换；hashLock 是两侧锁定的公共密码学锚点；timeLock 是失败退出机制。三者共同构成跨链状态机的关键上下文。


上述流程描述的是合约链（HTLC）路径。在无脚本链的 adaptor signature 路径下，总体角色协作和阶段划分保持一致，但具体的条件锁定方式有所不同：资产锁定不通过链上 HTLC 脚本表达，而是通过预签名（PreSign）绑定秘密点 T；资产释放不通过链上提交 secret 触发，而是通过持有秘密 t 的一方补全签名并广播交易完成；秘密传递不通过链上 secret 公开，而是另一方从链上出现的有效签名与预签名的差值中提取 t。失败退出机制则通过预设的 timeLock 退款交易或多签回退路径实现。两种路径在架构层共享同一套角色模型和协调流程，差异仅在于底层条件锁定的实现层——这正是系统能够同时覆盖智能合约链和无脚本链的关键。
